%---------------------------------------------------------------
\chapter{\babLima}
%\begin{center}
%	\textbf{BAB IV}
%	\par \textbf{ANALISIS SIMULASI SISTEM}
%\end{center}

\section{Simpulan}

Penelitian ini bertujuan untuk mengembangkan sistem rekomendasi \textit{Faculty Supervisor} berbasis \textit{semantic similarity} sebagai pendukung keputusan dalam proses pemetaan mahasiswa peserta Program Enrichment di XYZ University. Berdasarkan hasil analisis, perancangan, implementasi, dan evaluasi yang telah dilakukan, dapat ditarik simpulan yang disusun untuk menjawab ketiga rumusan masalah penelitian sebagai berikut.

\begin{enumerate}
    \item Menjawab rumusan masalah pertama, sistem rekomendasi \textit{Faculty Supervisor} magang berbasis website telah berhasil dirancang dan dikembangkan menggunakan pendekatan \textit{semantic similarity} dengan arsitektur \textit{Transformer} (model BAAI/bge-m3). Sistem ini menggabungkan skor kemiripan semantik, aturan \textit{company group bonus}, dan algoritma \textit{greedy capacity solver} untuk menjaga keseimbangan kuota pembimbingan setiap dosen.

    \item Menjawab rumusan masalah kedua, tingkat kesesuaian hasil rekomendasi sistem dengan \textit{ground truth} penugasan institusional tergolong tinggi. Pengujian menghasilkan \textit{Mean Reciprocal Rank} (MRR) sebesar 0,711, Hit@5 sebesar 0,929, dan \textit{Assignment Match Rate} sebesar 0,530, yang menunjukkan bahwa pembimbing relevan secara konsisten muncul pada peringkat teratas rekomendasi.

    \item Menjawab rumusan masalah ketiga, sistem yang dikembangkan terbukti meningkatkan efisiensi proses pemetaan \textit{Faculty Supervisor} magang oleh EPC. Waktu penentuan yang sebelumnya membutuhkan 1 hingga 3 minggu per \textit{batch} dapat dipersingkat menjadi hitungan menit, dengan proses yang lebih terotomatisasi, transparan, dan objektif, serta sebaran beban pembimbingan yang sangat merata antar dosen, ditandai dengan nilai \textit{Gini Coefficient} sebesar 0,0138.

\end{enumerate}

\section{Saran}

Berdasarkan hasil penelitian yang telah dilakukan, terdapat beberapa saran untuk pengembangan lebih lanjut:

\begin{enumerate}
    \item \textbf{Perluasan cakupan data supervisor.} Representasi dosen pembimbing saat ini bergantung pada data historis bimbingan dan rekomendasi topik skripsi yang tersedia. Penambahan sumber data seperti publikasi ilmiah, profil penelitian, atau portofolio proyek dosen dapat memperkaya representasi semantik dan meningkatkan kualitas rekomendasi.

    \item \textbf{Perluasan cakupan program enrichment.} Sistem yang dikembangkan saat ini difokuskan pada jalur magang (\textit{internship}). Pengembangan selanjutnya dapat mencakup jalur lain dalam Program Enrichment, seperti penelitian (\textit{research}), kewirausahaan (\textit{entrepreneurship}), dan pengembangan profesional, dengan penyesuaian pada skema representasi data mahasiswa.

    \item \textbf{\textit{Fine-tuning} model \textit{embedding} pada domain akademik Indonesia.} Model \textit{embedding} yang digunakan merupakan model \textit{pre-trained} bersifat umum. Melakukan \textit{fine-tuning} pada korpus akademik berbahasa Indonesia atau data spesifik domain Program Enrichment berpotensi meningkatkan akurasi pencocokan semantik secara signifikan.

    \item \textbf{Evaluasi usabilitas dengan pengguna nyata.} Pengujian yang dilakukan dalam penelitian ini terbatas pada \textit{black box testing} fungsional. Evaluasi lanjutan menggunakan \textit{System Usability Scale} (SUS) dengan melibatkan EPC sebagai pengguna nyata perlu dilakukan untuk mengukur tingkat penerimaan dan kemudahan penggunaan sistem secara empiris.

    \item \textbf{Integrasi dengan sistem informasi akademik institusi.} Saat ini sistem memerlukan unggah data manual dalam format Excel. Integrasi langsung dengan sistem informasi akademik XYZ University akan mengurangi potensi kesalahan \textit{input} data dan mempersingkat alur kerja EPC secara keseluruhan.

    \item \textbf{Studi longitudinal terhadap kualitas rekomendasi.} Evaluasi dalam penelitian ini menggunakan data historis satu tahun sebagai \textit{ground truth}. Studi lanjutan yang mengikuti keluaran rekomendasi selama beberapa periode pemetaan akan memberikan gambaran yang lebih representatif terhadap kualitas sistem dalam kondisi operasional nyata.

    \item \textbf{Replikasi sistem ke fakultas lain sebagai uji generalisabilitas.} Evaluasi penelitian ini terbatas pada 168 mahasiswa dan 14 supervisor dari satu program studi dalam satu batch pemetaan. Mengaplikasikan sistem pada fakultas lain dengan skala, komposisi supervisor, dan karakteristik data yang berbeda akan mengonfirmasi apakah model \textit{embedding} dan \textit{greedy solver} yang dikembangkan berperilaku konsisten di luar konteks Computer Science Program Enrichment.
\end{enumerate}
